\chapter{Eigenheiten}

\begin{rulebox}{Eigenheiten}
	Eigenheiten sind permanente Vor- und Nachteile.
	\begin{itemize}
		\item Jeder Charakter beginnt mit \textbf{1–2 Eigenheiten}.
		\item Eigenheiten skalieren nicht mit Rang.
		\item Sie können nur durch \textbf{Spielereignisse} erlangt oder verloren werden.
		\item Effekte sind dauerhaft aktiv, sofern nicht anders angegeben.
	\end{itemize}
\end{rulebox}

Eigenheiten sind thematisch in Gruppen unterteilt.  
Diese Gruppen dienen der Übersicht, Begrenzung bei der Charaktererschaffung
und der gezielten Erweiterung im Spiel.



	\begin{itemize}
		\item \textbf{Eigenheiten – Magie und Zauberei}  
		(z.\,B. instabile Magie, arkane Sensibilität, Manasucht)
		
		\item \textbf{Eigenheiten – Geist und Psyche}  
		(z.\,B. Phobien, Zwangsvorstellungen, eiserne Disziplin)
		
		\item \textbf{Eigenheiten – Kampf und Taktik}  
		(z.\,B. Überaggressiv, defensiver Instinkt, Tunnelblick)
		
		\item \textbf{Eigenheiten – Übernatürliches und Flüche}  
		(z.\,B. verfluchte Aura, Stimmen aus dem Jenseits, Schattenmal)
		
		\item \textbf{Eigenheiten – Herkunft und Umgebung}  
		(z.\,B. Straßenkind, Adelsblut, Wildnisgeboren)
		
		\item \textbf{Eigenheiten – Sozial und Ruf}  
		(z.\,B. berüchtigt, Ehrenkodex, Schuldverpflichtung)
		
		\item \textbf{Eigenheiten – Fähigkeiten und Fokus}  
		(z.\,B. einseitige Ausbildung, Perfektionist, Improvisator)
	\end{itemize}

\section{Eigenheiten Geist und Psyche}

\begin{traitbox}{\term{Phobie}{eigenheit-phobie}}
	
	\textbf{Nachteil:}  
	Der Charakter leidet unter einer tief verwurzelten Angst
	(z.\,B. Dunkelheit, bestimmte Kreaturen oder Magie; Wahl bei Erwerb).
	Bei Konfrontation ist eine \textbf{GESELLSCHAFT-Probe} erforderlich.
	Misslingt diese, \textbf{flieht} der Charakter und kann an der
	\textbf{gesamten Szene nicht teilnehmen}.
	
	\medskip
	\textbf{Vorteil:}  
	Aufgrund permanenter Anspannung und Wachsamkeit erhält der Charakter
	einen \textbf{+2 Bonus auf WAHRNEHMUNG}.
	
\end{traitbox}

\begin{traitbox}{\term{Wahnsinniges Genie}{eigenheit-wahnsinniges-genie}}
	
	\textbf{Nachteil:}  
	Der Charakter denkt obsessiv und verliert sich in Gedankenspiralen.
	Bei \textbf{wichtigen Entscheidungen} muss eine
	\textbf{WAHRNEHMUNG-Probe} abgelegt werden.
	Bei Misserfolg handelt der Charakter \textbf{irrational oder unlogisch}
	(Interpretation durch die Spielleitung).
	
	\medskip
	\textbf{Vorteil:}  
	Der Charakter erhält:
	\begin{itemize}
		\item \textbf{+2 Bonus auf WISSEN-Proben}
		\item Einen frei wählbaren \textbf{Basis-Skill auf Rang 2}  
		(\textbf{Magierklassen} dürfen stattdessen einen \textbf{Basiszauber} wählen)
	\end{itemize}
	
\end{traitbox}

\begin{traitbox}{\term{Fluch des Wissens}{eigenheit-fluch-des-wissens}}
	
	\medskip
	\textbf{Nachteil:}
	Eine übernatürliche Wissensquelle plagt den Charakter mit Visionen. Nach jeder WISSEN-Probe besteht die Chance, für 1W4 Runden handlungsunfähig zu werden. 

	\medskip
	\textbf{Vorteil:}
	Einmal pro Szene kann eine Frage zu einem Mysterium gestellt werden – die Antwort erfolgt kryptisch durch die Quelle (Spielleiter).
	
\end{traitbox}

\section{Eigenheiten Magie und Zauberei}

\begin{traitbox}{\term{Vergesslichkeit}{eigenheit-vergesslichkeit}}
	
	\textbf{Nachteil:}  
	Bei jedem Wirken eines Zaubers würfelt der Charakter einen \textbf{W6}.
	Bei einer \textbf{1} schlägt der Zauber fehl und wird nicht ausgeführt
	(Kosten fallen regulär an).
	
	\medskip
	\textbf{Vorteil:}  
	Der Charakter erhält \textbf{+2 zusätzliche Zauberpunkte}
	beim Erstellen oder Steigern des Charakters.
	
\end{traitbox}

\begin{traitbox}{\term{Magische Fehlzündung}{eigenheit-magische-fehlzuendung}}
	
	\textbf{Nachteil:}  
	Bei jedem Zauberwirken besteht eine \textbf{1W4-Chance} auf eine Fehlzündung.
	Bei einer gewürfelten \textbf{1} tritt ein \textbf{unkontrollierter magischer Nebeneffekt}
	auf (Art und Auswirkung entscheidet die Spielleitung).
	
	\medskip
	\textbf{Vorteil:}  
	Der Charakter besitzt eine \textbf{wildmagische Resonanz}.  
	Erzielt er beim Zauberwirken einen \textbf{kritischen Erfolg (Ass-Mechanik)},
	wird der Effekt verstärkt (z.\,B. erhöhte Dauer, Reichweite oder Schaden;
	konkret durch die Spielleitung festgelegt).
	
\end{traitbox}

\begin{traitbox}{\term{Manaschwäche}{eigenheit-manaschwaeche}}
	
	\textbf{Nachteil:}  
	\begin{itemize}
		\item \textbf{–2} auf nächtliche Manaregeneration
		\item \textbf{–2 Zauberpunkte} zu Spielbeginn
	\end{itemize}
	
	\medskip
	\textbf{Vorteil:}  
	Wähle \textbf{eine} der folgenden Optionen (bei Erwerb festlegen):
	\begin{itemize}
		\item \textbf{1× pro Tag} einen bekannten Zauber ohne Manakosten wirken
		\item Alle Zauber kosten \textbf{–2 Manapunkte (MP)}
		(\textbf{SL-Zustimmung erforderlich}, insbesondere bei hochrangigen Zaubern)
	\end{itemize}
	
\end{traitbox}

\begin{traitbox}{\term{Gier nach Macht}{eigenheit-gier-nach-macht}}
	
	\textbf{Nachteil:}  
	Der Charakter ist zwanghaft getrieben, jede Quelle magischer Macht zu untersuchen.
	Dies kann ihn regelmäßig in \textbf{gefährliche oder unkontrollierbare Situationen}
	bringen (Auslegung durch die Spielleitung).
	
	\medskip
	\textbf{Vorteil:}  
	Der Charakter besitzt einen natürlichen magischen Spürsinn:
	Er erhält einen \textbf{+1 Bonus} auf alle Proben zur Analyse magischer
	Gegenstände oder Orte
	(z.\,B. \textbf{WISSEN}, \textbf{Magiesinn}, \textbf{Arkaner Spürsinn}).
	
\end{traitbox}

\begin{traitbox}{\term{Verfluchte Gabe}{eigenheit-verfluchte-gabe}}
	
	\textbf{Nachteil:}  
	Eine fremde Macht fordert regelmäßig Tribut.
	Nach jedem Abenteuer oder einer längeren Szene kann eine
	\textbf{unerwartete Nebenwirkung} auftreten
	(z.\,B. Albträume, ein spontaner Zaubereffekt oder fremdartige Stimmen;
	Art entscheidet die Spielleitung).
	
	\medskip
	\textbf{Vorteil:}  
	\textbf{1× pro Sitzung} kann der Charakter einen \textbf{mächtigen Zauber}
	wirken, den er normalerweise nicht beherrscht.
	Art und Herkunft des Effekts werden vorab mit der Spielleitung festgelegt.
	
\end{traitbox}

\section{Eigenheiten Kampf und Taktik}

\begin{traitbox}{\term{Im Kampf verkrüppelt}{eigenheit-im-kampf-verkrueppelt}}
	
	\textbf{Nachteil:}  
	Körperliche Einschränkungen verursachen einen
	\textbf{–1 Malus} auf alle \textbf{STÄRKE}- und \textbf{GEWANDTHEIT}-Proben.
	
	\medskip
	\textbf{Vorteil:}  
	Durch taktisches Denken erhält der Charakter einen
	\textbf{+1 Bonus} auf alle \textbf{KAMPF}- und \textbf{WISSEN}-Proben.
	
\end{traitbox}

\begin{traitbox}{\term{Ritterlichkeit}{eigenheit-ritterlichkeit}}
	
	\textbf{Nachteil:}  
	Der Charakter unterliegt einem strengen Ritterkodex.
	Die Spielleitung führt einen verdeckten \textbf{Ehrezähler} (Startwert: 10).
	Verstöße senken den Wert und können situationsabhängige Mali oder Konsequenzen auslösen. Am Beginn jedes Abenteuers erhält der Spieler den aktuellen Ehrezähler. 
	
	\medskip
	\textbf{Vorteil:}  
	Der Charakter ist \textbf{immun gegen Statuseffekte}.
	
	\medskip
	\textbf{Ritterkodex (bindend):}
	\begin{itemize}
		\item Beschütze die Schwachen
		\item Bewahre deine Ehre (kein Wortbruch, keine Heimtücke, kein Verrat)
		\item Achte die Gesetze des Krieges (Zivilisten, religiöse Stätten)
		\item Führe ehrenvolle Duelle (keine Hinterhalte, Niederlagen akzeptieren)
	\end{itemize}
	
	\medskip
	\textbf{Hinweis für den SL:}
	Immunität gegen Statuseffekte ist vergleichsweise recht stark, deshalb sollte das führen des Ehrezählers strikt durchgeführt werden und entsprechende Konsequenzen mit unter herausfordernd (aber nicht unfair!) sein.
\end{traitbox}

\begin{traitbox}{\term{Kampf in Trance}{eigenheit-kampf-in-trance}}
	
	\textbf{Nachteil:}  
	Im Kampfrausch muss der Charakter eine \textbf{WAHRNEHMUNG}-Probe ablegen,
	um Freund von Feind zu unterscheiden.
	Bei Misserfolg entscheidet die Spielleitung über das Ziel der Aktion.
	
	\medskip
	\textbf{Vorteil:}  
	Der Charakter erhält einen \textbf{+2 Bonus auf alle Angriffsproben}
	während des Kampfes.
	
\end{traitbox}

\begin{traitbox}{\term{Meister der Improvisation}{eigenheit-meister-der-improvisation}}
	
	\textbf{Nachteil:}  
	\textbf{–1 Malus} auf den Einsatz klassischer oder spezialisierter Waffen.
	
	\medskip
	\textbf{Vorteil:}  
	\textbf{+2 Bonus} auf kreative Einsätze von
	Umgebung, Alltagsgegenständen oder improvisierten Waffen.
	
\end{traitbox}

\section{Eigenheiten Übernatürliches und Flüche}

\begin{traitbox}{\term{Preis des Lebens}{eigenheit-preis-des-lebens}}
	
	\textbf{Nachteil:}  
	Die Rückkehr aus dem Tod hat Risse hinterlassen.
	In emotionalen Extremsituationen (große Angst, starker Schmerz, Nähe zum Tod)
	erleidet der Charakter \textbf{1W6 Seelenschaden} durch einen \textbf{Schattenriss}
	(Auslösung nach SL-Entscheidung).
	
	\medskip
	\textbf{Vorteil:}  
	Der Charakter steht mit einem Fuß außerhalb der Zeit.
	\textbf{1× pro Tag} darf er eine \textbf{soeben misslungene Probe}
	\textbf{zurückdrehen und erneut würfeln}.
	Der neue Wurf erhält \textbf{+1 Bonus}.
	
	\medskip
	\textbf{SL-Hinweis:}  
	Während eines Schattenrisses kann der Charakter Stimmen hören,
	fremde Erinnerungen erleben oder von Kreaturen der Dunkelheit wahrgenommen/angegriffen werden (z.B. Silithar).
	
\end{traitbox}

\begin{traitbox}{\term{Schattenfreund}{eigenheit-schattenfreund}}
	
	\textbf{Nachteil:}  
	Tiere und einfache Menschen meiden den Charakter instinktiv.
	\textbf{–2 Malus} auf soziale Interaktionen mit Tieren und einfachen Bürgern.
	
	\medskip
	\textbf{Vorteil:}  
	Schatten scheinen ihn zu schützen.
	In dämmriger oder dunkler Umgebung erhält der Charakter
	\textbf{+2 Bonus auf HEIMLICHKEIT}.
	Zusätzlich kann er sich \textbf{1× pro Tag} für \textbf{1 Runde unsichtbar} machen.
	
\end{traitbox}

\begin{traitbox}{\term{Fluch der Vorfahren}{eigenheit-fluch-der-vorfahren}}
	
	\textbf{Nachteil:}  
	Eine uralte Schuld lastet auf dem Charakter.
	In zivilisierten Regionen erhält er einen
	\textbf{–1 Malus auf GESELLSCHAFT} durch Misstrauen und alte Gerüchte.
	
	\medskip
	\textbf{Vorteil:}  
	Der Fluch entfaltet Macht in Grenzsituationen.
	Fällt der Charakter auf \textbf{5 oder weniger Lebenspunkte (LP)},
	erhält er für \textbf{2 Runden einen +2 Bonus auf alle Proben}.
	
\end{traitbox}

\begin{traitbox}{\term{Dämonischer Begleiter}{eigenheit-daemonischer-begleiter}}
	
	\textbf{Nachteil:}  
	Der Charakter wird von einem Dämon begleitet,
	der ihn manipuliert, provoziert oder zu riskanten Entscheidungen drängt
	(Interpretation und Einfluss nach SL-Entscheidung).
	
	\medskip
	\textbf{Vorteil:}  
	Der Dämon kann den Charakter aktiv unterstützen,
	z.\,B. durch Hinweise, Warnungen oder das Auslösen eines
	kleinen magischen Effekts (Häufigkeit und Stärke nach SL-Freigabe).
	
\end{traitbox}

\begin{traitbox}{\term{Zweites Gesicht}{eigenheit-zweites-gesicht}}
	
	\textbf{Nachteil:}  
	Der Charakter sieht regelmäßig Dinge, die andere nicht wahrnehmen
	(Geister, Visionen, Ebenenfragmente).
	Dies führt zu Ablenkung:
	\textbf{–1 Malus auf WAHRNEHMUNG} in Alltagssituationen
	und gelegentlich \textbf{Verwirrung} (SL-Entscheidung).
	
	\medskip
	\textbf{Vorteil:}  
	Der Charakter erkennt Lügen, Illusionen und magische Einflüsse intuitiv.
	\textbf{+2 Bonus} auf entsprechende Proben.
	Zusätzlich kann er \textbf{1× pro Sitzung ein Omen erfragen}
	(kryptische Warnung oder Hilfe durch die SL).
	
\end{traitbox}

\begin{traitbox}{\term{Geisterflüstern}{eigenheit-geisterfluester}}
	
	\textbf{Nachteil:}  
	Stimmen aus der Geisterwelt begleiten den Charakter ständig.
	\textbf{–2 Malus} auf alle \textbf{GESELLSCHAFT}- und \textbf{WAHRNEHMUNG}-Proben.
	
	\medskip
	\textbf{Vorteil:}  
	Der Charakter kann \textbf{Geisterwesen wahrnehmen und mit ihnen kommunizieren}.
	Je nach Situation können Geister Hinweise geben,
	Informationen liefern oder in kritischen Momenten indirekt helfen
	(Art und Umfang nach SL-Entscheidung).
	
\end{traitbox}

\section{Eigenheiten Herkunft und Umgebung}

\begin{traitbox}{\term{In der Natur aufgewachsen}{eigenheit-natur-aufgewachsen}}
	
	\textbf{Nachteil:}  
	Der Charakter kennt kein Eigentum und greift instinktiv zu.
	\textbf{–2 Malus auf GESELLSCHAFT}.
	Bei Diebstahl besteht nur eine \textbf{1W4-Chance}, unentdeckt zu bleiben.
	
	\medskip
	\textbf{Vorteil:}  
	Tiefe Verbundenheit mit der Wildnis.
	\textbf{+2 Bonus auf WILDNISLEBEN}.
	
\end{traitbox}

\begin{traitbox}{\term{Im Feuer geboren}{eigenheit-feuer-geboren}}
	
	\textbf{Nachteil:}  
	Der Charakter reagiert extrem empfindlich auf Kälte.
	\textbf{Kälteschaden und Unterkühlung wirken zu 150\%}.
	
	\medskip
	\textbf{Vorteil:}  
	Der Charakter ist \textbf{immun gegen Feuerschaden und Verbrennung}.
	
\end{traitbox}

\begin{traitbox}{\term{In Dunkelheit geboren}{eigenheit-dunkelheit-geboren}}
	
	\textbf{Nachteil:}  
	Bei hellem Tageslicht erhält der Charakter einen
	\textbf{–1 Malus auf WAHRNEHMUNG}.
	
	\medskip
	\textbf{Vorteil:}  
	An die Nacht angepasst.
	Bei Dunkelheit erhält der Charakter
	\textbf{+1 Bonus auf WAHRNEHMUNG}.
	Nachtsicht ist möglich (Reichweite und Details nach SL-Festlegung).
	
\end{traitbox}

\begin{traitbox}{\term{Alchemistischer Körper}{eigenheit-alchemistischer-koerper}}
	
	\textbf{Nachteil:}  
	Der Körper reagiert instabil auf Heilmittel.
	Bei jeder Tranknutzung besteht eine \textbf{1W4-Chance},
	dass eine unerwartete Nebenwirkung eintritt
	(Art nach SL-Entscheidung).
	
	\medskip
	\textbf{Vorteil:}  
	Der Charakter ist \textbf{immun gegen Giftschaden und Vergiftung}.
	Zusätzlich erhält er
	\textbf{+2 Bonus auf WISSEN-Proben} in
	Alchemie, Anatomie und Giftkunde.
	
\end{traitbox}

\section{Eigenheiten Sozial und Ruf}

\begin{traitbox}{\term{Eidbrecher}{eigenheit-eidbrecher}}
	
	\textbf{Nachteil:}  
	Der Charakter hat in der Vergangenheit einen Eid gebrochen.
	Er erleidet einen \textbf{–2 Malus auf GESELLSCHAFT}
	bei Interaktionen mit ehrenhaften, religiösen oder schwurgebundenen Gruppen.
	
	\medskip
	\textbf{Vorteil:}  
	Schuld und Zweifel haben keine Macht mehr über ihn.
	Der Charakter ist \textbf{resistent gegen psychologische Konsequenzen}
	(z.\,B. Schuldgefühle, Reue, Gewissenskonflikte; SL bestimmt konkrete Effekte).
	
\end{traitbox}

\begin{traitbox}{\term{Schlechter Ruf}{eigenheit-schlechter-ruf}}
	
	\textbf{Nachteil:}  
	Der Charakter ist in bestimmten Regionen oder Kreisen bekannt und unbeliebt.
	Er erhält dort einen \textbf{–2 Malus auf GESELLSCHAFT}.
	
	\medskip
	\textbf{Vorteil:}  
	Der Charakter verfügt über zwielichtige Kontakte.
	Er kann Informationen, Unterschlupf oder Dienstleistungen erhalten,
	wenn diese glaubwürdig aus dem Milieu stammen
	(Art und Umfang nach SL-Entscheidung).
	
\end{traitbox}

\begin{traitbox}{\term{Unheimliche Präsenz}{eigenheit-unheimliche-praesenz}}
	
	\textbf{Nachteil:}  
	Andere fühlen sich unwohl in seiner Nähe.
	Der Charakter erhält einen \textbf{–2 Malus auf GESELLSCHAFT}
	bei positiven sozialen Interaktionen
	(z.\,B. Überzeugen, Freundschaften, Diplomatie).
	
	\medskip
	\textbf{Vorteil:}  
	Seine Aura wirkt einschüchternd.
	Der Charakter erhält einen \textbf{+2 Bonus auf GESELLSCHAFT}
	bei negativen sozialen Interaktionen
	(z.\,B. Drohungen, Erpressung, offene Manipulation).
	
\end{traitbox}

\begin{traitbox}{\term{Pazifist}{eigenheit-pazifist}}
	
	\textbf{Nachteil:}  
	Der Charakter verweigert tödliche Gewalt.
	Er kann keine Angriffe ausführen, die tödlich enden könnten
	(z.\,B. Waffenangriffe gegen Humanoide, tödliche Zauber).
	In Gruppenkämpfen muss er eine \textbf{WAHRNEHMUNG-Probe} bestehen,
	um sich überhaupt an gewalttätigen Handlungen zu beteiligen. Unterstützende Handlungen (Heilung, Schutz, Ablenkung)
	sind weiterhin erlaubt.
	
	
	\medskip
	\textbf{Vorteil:}  
	Der Charakter strahlt Frieden und Entschlossenheit aus.
	Er erhält einen \textbf{+3 Bonus auf GESELLSCHAFT} bei Verhandlungen.
	Zusätzlich kann er \textbf{1× pro Spielsitzung}
	einen Gegner mit einer Probe auf \textbf{GESELLSCHAFT + WISSEN}
	zum Aufgeben oder Innehalten bewegen
	(\textbf{nicht} gegen Bosse oder erzählerisch zentrale Gegner).
	Feindliche Kreaturen ignorieren ihn häufig,
	solange er nicht angreift (SL-Entscheidung).
	
\end{traitbox}

\begin{traitbox}{\term{Alter Geldadel}{eigenheit-geldadel}}
	
	\textbf{Nachteil:}  
	Der Charakter ist es gewohnt, Probleme mit Geld zu lösen.
	Er erhält einen \textbf{–2 Malus auf WILDNISLEBEN und HEIMLICHKEIT}
	bei Improvisation, Reisen oder einfachen Tätigkeiten.
	
	\medskip
	\textbf{Vorteil:}  
	Der Charakter beginnt mit \textbf{deutlich erhöhtem Startkapital}
	(empfohlen: dreifaches Startvermögen).
	Zusätzlich kann er \textbf{1× pro Spielsitzung}
	eine Zahlung von \emph{hohem, aber glaubwürdigem Betrag} leisten,
	um ein Problem zu lösen
	(z.\,B. Bestechung, Sonderausrüstung, Söldnerhilfe).
	
\end{traitbox}

\begin{traitbox}{\term{Verarmt}{eigenheit-verarmt}}
	
	\textbf{Nachteil:}  
	Der Charakter ist ungeübt im Umgang mit Macht, Reichtum und Autoritäten.
	Er erhält einen \textbf{–2 Malus auf GESELLSCHAFT}
	in nobler Umgebung oder bei Gesprächen mit Obrigkeit.
	
	\medskip
	\textbf{Vorteil:}  
	Der Charakter kennt den Überlebenskampf.
	Er erhält einen \textbf{+2 Bonus auf HEIMLICHKEIT und WILDNISLEBEN}
	in städtischem, verwahrlostem oder unwirtlichem Gelände
	(z.\,B. Gassen, Ruinen, Wildnis).
	
\end{traitbox}

\begin{traitbox}{\term{Nutzlos}{eigenheit-nutzlos}}
	
	\textbf{Nachteil:}  
	Der Charakter wird unterschätzt oder ignoriert.
	Er erhält einen \textbf{–2 Malus auf GESELLSCHAFT},
	wenn er sich behaupten oder ernst genommen werden möchte.
	
	\medskip
	\textbf{Vorteil:}  
	Gerade weil man ihn für irrelevant hält, entgeht er oft der Aufmerksamkeit.
	Er erhält einen \textbf{+2 Bonus auf HEIMLICHKEIT}
	in brenzligen Situationen.
	Gegner übersehen ihn häufig in der ersten Kampfrunde
	(SL-Entscheidung).
	
\end{traitbox}

\section{Eigenheiten Fähigkeiten und Fokus}

\begin{traitbox}{\term{Geniales Talent}{eigenheit-geniales-talent}}
	
	\textbf{Nachteil:}  
	Der Charakter ist extrem auf ein einzelnes Talent fokussiert.
	Alle anderen Talente erhalten dauerhaft \textbf{–1 Talentpunkt}
	(ausgenommen das gewählte Talent).
	
	\medskip
	\textbf{Vorteil:}  
	Der Charakter erhält \textbf{+2 Talentpunkte}
	auf ein frei wählbares Talent.
	
\end{traitbox}

\begin{traitbox}{\term{Vergessener Meister}{eigenheit-vergessener-meister}}
	
	\textbf{Nachteil:}  
	Der Charakter leidet unter wiederkehrenden Gedächtnislücken.
	Zu Beginn jedes Abenteuers (Referenz: 1 Abenteuer = 1-2 Abende Spielzeit/ eine Geschichte in der Kampagne)
	verliert er \textbf{1–2 Zauberpunkte oder Skillpunkte}.
	Für jeden Punkt darf eine \textbf{1W6-Probe} abgelegt werden;
	bei Erfolg wird der jeweilige Punkt nicht verloren.
	
	\medskip
	\textbf{Vorteil:}  
	Der Charakter startet mit \textbf{+5 zusätzlichen Zauberpunkten oder Skillpunkten}
	(Spielerentscheidung bei Erstellung).
	
\end{traitbox}

\begin{traitbox}{\term{Gebrechlich durchs Alter}{eigenheit-gebrechlich-alter}}
	
	\textbf{Nachteil:}  
	Der Körper zeigt deutliche Alterserscheinungen.
	Der Charakter erhält:
	\begin{itemize}
		\item \textbf{–4 Lebenspunkte (LP)}
		\item \textbf{–1 Malus auf STÄRKE}
		\item \textbf{–1 Malus auf GEWANDTHEIT}
	\end{itemize}
	
	\medskip
	\textbf{Vorteil:}  
	Lebenserfahrung gleicht körperliche Schwäche aus.
	Der Charakter erhält \textbf{+4 Zauberpunkte oder Skillpunkte}
	(frei wählbar).
	
\end{traitbox}

